% Imporditud: tools/import_eksperimendid.py
% Allikas: Eesti füüsikaolümpiaadi eksperimendi ülesannete
%   kogu (Rael Kalda, 2023), ülesanne Ü90, lahendus L90.
\setAuthor{EFO žürii}
\setRound{piirkonnavoor}
\setYear{2003}
\setNumber{G E2}
\setDifficulty{3}
\setTopic{geomeetriline optika}
\setVahendid{koondav lääts, laelamp}
\setPunktid{12}
\setAllikas{Eksperimendi kogu Ü90, lahendus lk 72}
\setLahendusseis{toores}

\prob{Kaks kujutist}
Kui vaadata kumerläätselt peegeldunud valgust, näeme valgusallika kahte kujutist. Miks? Mille poolest need kujutised erinevad? Kumb neist asub meile lähemal? Kuidas seda eksperimentaalselt kontrollida?

\hint

\solu
a) Kujutised tekivad valguse peegeldumisel läätse esimeselt ja tagumiselt pinnalt (1 p). b) Läätse tagumiselt, nõgusalt pinnalt peegeldumisel tekkiva kujutise konstrueerimine (2 p).

c) See on tõeline kujutis (1 p). d) Tõeline kujutis on ümberpööratud (1 p). e) Läätse eesmiselt, kumeralt pinnalt pinnalt peegeldumisel tekkiva kujutise konstrueerimine (2 p). f) See on näiv kujutis (1 p). g) Näiv kujutis on päripidine (1 p). h) Mõlemad kujutised on vähendatud, kuid tõeline kujutis on väiksem (1 p), sest esimeselt pinnalt toimub ainult peegeldumine, tagumiselt aga lisaks veel murdumine läätses: see lühendab fookuskaugust (ilma murdumiseta oleksid need mõlemal juhul võrdsed) ja seepärast on kujutis väiksem. i) Meile lähemal asub tõeline kujutis (1 p). j) Kauguse kindlakstegemiseks võib näiteks silma lähendada läätsele. Kui kaugus kujutisest saab väga väikeseks, muutub kujutis ebateravaks (1 p) Kaugusi saab kindlaks teha ka parallaksi meetodil, liigutades pead ja leides juurdeviidud eseme (pliiatsi) sellise asendi, kus see ese kujutise suhtes ei liigu.
\probend
